\subsection{A conductor class retained by supported singularity theory}

\begin{lemma}[Proposed: the determinant of a localized perfect object]
\label{C:supported-singularity-Picard-detector}
Let $D$ be a commutative noetherian local ring and let $h\in D$.
Write $D_h=D[1/h]$ and let
\[
 \mathcal T_h(D)=
 \ker\bigl(D_{\rm sg}(D)\longrightarrow D_{\rm sg}(D_h)\bigr)
                                                               \tag{C.1.PP1}
\]
be the full triangulated kernel, without idempotent completion.
There is a natural surjection
\[
 K_0(\mathcal T_h(D))\longrightarrow\operatorname{Pic}(D_h),
 \qquad [X]\longmapsto [\det(X\otimes_D D_h)].        \tag{C.1.PP2}
\]
In particular the Grothendieck group of this supported category
need not inject into the Grothendieck group of the ambient
singularity category.
\end{lemma}
\begin{proof}
Let $\mathcal C_h\subset D^b(\operatorname{mod}D)$ be the full
subcategory of complexes whose localization is perfect over $D_h$.
It contains $\operatorname{Perf}(D)$.
The quotient $\mathcal C_h/\operatorname{Perf}(D)$ identifies
with the full subcategory (C.1.PP1): its objects are exactly those
whose localizations become zero in $D_{\rm sg}(D_h)$.
Full faithfulness follows directly from roofs in the Verdier
quotient.  If one endpoint of a denominator belongs to
$\mathcal C_h$ and its cone is perfect over $D$, the other endpoint
also belongs to $\mathcal C_h$.  Thus all the roofs and their
equivalences can already be taken inside this subcategory.

For a bounded complex $P^\bullet$ of finite projective $D_h$-modules,
define its determinant line to be the tensor product of the top
exterior powers of its terms, with odd terms dualized.
A bounded acyclic complex of projectives is split exact, so its
determinant is trivial.  Cone triangles give multiplication of
determinant classes.  Therefore localization and determinant
give an additive map from $K_0(\mathcal C_h)$ to
$\operatorname{Pic}(D_h)$.
Every perfect complex over the local ring $D$ is represented by
a bounded finite free complex, so this map vanishes on
$K_0(\operatorname{Perf}(D))$.
The generators and triangle relations for the Verdier quotient
then give (C.1.PP2).  Equivalently, roofs with perfect cones change
the determinant by a trivial line.

Finally, every line bundle $L$ over $D_h$ extends to a finite
$D$-module.  Indeed choose a finite presentation of $L$ and clear
the finitely many powers of $h$ in its presentation matrix; the
cokernel over $D$ localizes to $L$.
This module is in $\mathcal C_h$, and its determinant after
localization is $L$ itself.  The map is surjective.
The explicit noninjectivity assertion is supplied by the next
proposition.
\end{proof}

\begin{proposition}[Proposed: a nonzero supported class in the ramified cubic pencil]
\label{C:ramified-pencil-supported-Picard}
Let $A$ be a discrete valuation ring of mixed
characteristic $(0,2)$, with uniformizer $\pi$, residue field
$\kappa$, and $2\in\pi^2A$.
Set $b=2/\pi^2\in A$, choose $c\in A^\times$, and choose
either $D_0=A[h]_{(\pi,h)}$, or $D_0=A[[h]]$ with $A$
complete in the latter case.  Put
\[
 \begin{split}
 f&=U^2+2U-(c+h)\pi^3,\\
 D&=D_0[U]/(f),\\
 D'&=D_0[W]/
       \bigl(W^2+\pi bW-\pi(c+h)\bigr).
 \end{split}                                                   \tag{C.1.PP3}
\]
Write $C_0=D_0/\pi D_0$ and $F_0=C_0[1/h]$.
Thus $(C_0,F_0)$ is $(\kappa[h]_{(h)},\kappa(h))$ in the
polynomial-local case and $(\kappa[[h]],\kappa((h)))$ in
the complete formal case.
The map $D\hookrightarrow D'$ given by $U=\pi W$ identifies
$D'$ with the finite normalization of $D$.
Its conductor is $I=\pi D'=(\pi,U)D$, and
\[
 D/I=C_0,\qquad D'/I=C_0[W]/(W^2).                  \tag{C.1.PP4}
\]
There is an isomorphism of abelian groups
\[
 \operatorname{Pic}(D_h)
       \simeq F_0/C_0,                              \tag{C.1.PP5}
\]
where the clutching unit $1+W\alpha$ corresponds to the class
of $\alpha\in F_0$.

The explicit maximal Cohen--Macaulay $D$-module
\[
                  M=\pi D_0+(h+W)D_0\subset D'       \tag{C.1.PP6}
\]
belongs to $\mathcal T_h(D)$ and satisfies
\[
 [M]\ne0\quad\hbox{in }K_0(\mathcal T_h(D)),
 \qquad [M]=0\quad\hbox{in }K_0(D_{\rm sg}(D)).       \tag{C.1.PP7}
\]
Its image under (C.1.PP2), using (C.1.PP5), is
$h^{-1}+C_0$.
In its $D_0$-basis $(\pi,h+W)$, multiplication by $U$ is
\[
 X=\begin{pmatrix}
 -\pi h&\pi(c+h)-h^2+\pi b h\\
 \pi^2&\pi h-\pi^2b
 \end{pmatrix}.                                      \tag{C.1.PP8}
\]
Consequently an actual matrix factorization representing $M$ is
\[
       (UI-X,\ (U+2)I+X)                             \tag{C.1.PP9}
\]
over $D_0[U]$ with potential $f$, and over its localization
at $(\pi,h,U)$; in the complete case it also gives the
same matrix factorization over $A[[h,U]]$.
These assertions concern the supported pencil category.
They do not yet give a class in the Swan group of
$A[[x,z]]/(xz-\pi^3)$.
\end{proposition}
\begin{proof}
Each monic quadratic in (C.1.PP3) makes the displayed quotient
a free rank-two $D_0$-module.
In the formal case the assumption that $A$ is complete also
identifies these rings with the indicated formal-power-series
quotients: the finite $D_0$-algebra is complete, the additional
variable is topologically nilpotent, and completion does not
change it.  This identification is not needed in the
polynomial-local case, whose $D$ is exactly the local affine
pencil ring at $(\pi,h,U)$.
Thus
\[
             D=D_0\oplus\pi D_0W,
             \qquad D'=D_0\oplus D_0W.              \tag{C.1.PP10}
\]
The substitution is injective, and adjoining $W=U/\pi$ is
finite and birational.

In $D'$ the relation is
\[
                    W^2=\pi(c+h-bW).                \tag{C.1.PP11}
\]
The last factor is a unit.  The ring is local of dimension two,
with maximal ideal generated by $h,W$, because (C.1.PP11)
expresses $\pi$ as a multiple of $W^2$.
It is therefore regular, hence a normal domain and a UFD
\cite[Tag 0AG0]{C-NormalNodesStacks}.
It follows that $D$ is a domain and $D'$ its normalization:
an element of their common fraction field integral over $D$
is integral over $D'$ and hence belongs to the normal ring $D'$.
Also $D'/hD'$ is a one-dimensional regular local domain with
uniformizer $W$, or directly the Eisenstein quadratic extension
defined by $W^2+\pi bW-\pi c$.
In particular $h$ is a prime element of $D'$.

For $a+\pi dW\in D$, multiplication by $W$ gives
\[
 (a+\pi dW)W
       =\pi^2d(c+h)+(a-\pi^2bd)W.
\]
By (C.1.PP10) this belongs to $D$ exactly when $a\in\pi D_0$.
As $D'$ is generated by $1,W$ over $D_0$, this proves that
the conductor is precisely $\pi D'$.
Its two generators over $D$ are $\pi,U$, and reduction gives
(C.1.PP4).

Localize the resulting conductor square at $h$.
The Milnor patching theorem and its units--Picard exact sequence
\cite[I.2.6--2.7 and I.3.10]{C-WeibelKbook2013} give
\[
 (D'_h)^\times\times F_0^\times
   \longrightarrow
   \bigl(F_0[W]/W^2\bigr)^\times
   \longrightarrow\operatorname{Pic}(D_h)
   \longrightarrow0.                                \tag{C.1.PP12}
\]
Here $D'_h$ is a UFD and has trivial Picard group; the other
two quotient rings are a field and a local Artinian ring,
respectively, so their Picard groups are also zero.
The UFD implication follows, for example, from the injection
of the Picard group into the divisor class group
\cite[I.3.6--3.7]{C-WeibelKbook2013}.

Since $h$ is prime in the domain $D'$, every unit of $D'_h$
is $h^m u$ for $m\in\mathbb Z$ and $u\in(D')^\times$.
Indeed if two elements multiply to a power of $h$, repeated
primality of $h$ leaves units after the powers of $h$ are removed.
Modulo $I$ write
\[
             \bar u=\alpha(h)+W\beta(h),
 \qquad \alpha\in C_0^\times,
        \quad\beta\in C_0.
\]
After quotienting by the constant units $F_0^\times$,
the contribution of $h^m u$ is
$1+W\beta/\alpha$, with $\beta/\alpha\in C_0$.
Conversely every such element occurs: lift it through the
surjection $D_0\to C_0$ and use the unit
$1+W\widetilde\gamma(h)$ in $D'$.
All units in the middle group of (C.1.PP12) uniquely have the
form $a(1+W\gamma)$ with $a\in F_0^\times$ and
$\gamma\in F_0$.
Because $W^2=0$, multiplication adds the latter coordinates.
The exact quotient is therefore (C.1.PP5).

The two generators in (C.1.PP6) are $D_0$-independent.
Multiplication by $U=\pi W$, using (C.1.PP11), gives exactly
the two columns of (C.1.PP8), so $M$ is a $D$-module.
It is free of rank two over the regular ring $D_0$.
In particular $\pi,h$ is an $M$-regular sequence, so $M$ is
maximal Cohen--Macaulay over the dimension-two ring $D$.

The module $M$ contains $I=\pi D'$: it contains $\pi$, and
$\pi W=\pi(h+W)-h\pi$ also belongs to it.
Its image modulo $I$ is the submodule $(h+W)C_0$
of $C_0[W]/W^2$.
After inverting $h$, this is the same as
$(1+W/h)F_0$.
Consequently $M_h$ is the module patched from the free rank-one
modules over $D'_h$ and $F_0$ by the unit $1+W/h$.
Milnor patching makes it an actual line bundle on $D_h$.
Thus $M\in\mathcal T_h(D)$, and its determinant class is
$h^{-1}+C_0$, which is nonzero.
Lemma~\ref{C:supported-singularity-Picard-detector} proves the
first assertion of (C.1.PP7).

For the second assertion, the two maps
\[
 \begin{split}
 D'&\longrightarrow C_0,
             &a+dW&\longmapsto\bar d,\\
 D'&\longrightarrow C_0,
             &a+dW&\longmapsto\bar a-h\bar d
 \end{split}                                                   \tag{C.1.PP13}
\]
are surjective $D$-linear maps with kernels $D$ and $M$,
respectively.  Their $D$-linearity follows also by reduction
modulo $I$, since the image of $D$ is the constant subring
$C_0$ of (C.1.PP4).
They give two exact sequences with identical middle and final
terms, hence $[M]=[D]$ in $G_0(D)$.
Their triangles in $D_{\rm sg}(D)$ give $[M]=0$ there,
because $D$ is perfect.  This proves the claimed genuine
failure of injectivity for the full supported category.

Finally $\operatorname{tr}X=-\pi^2b=-2$ and
$\det X=-\pi^3(c+h)$, so Cayley--Hamilton gives
$X^2+2X=\pi^3(c+h)I$.
It follows that both products in (C.1.PP9) equal $fI$.
The usual polynomial-action presentation has cokernel $M$;
it remains exact after completion, or follows directly from
the displayed finite $D_0$-basis.
This verifies the actual matrix factorization.
\end{proof}

\noindent\emph{Precise remaining transfer.}
The ring in (C.1.PP3) is the quadratic part of the unit-curvature
chart $q-(c+h)(\pi^3-xz)$ for the node.
Identifying its $h$-supported category, after a relative
$xz$ suspension, with an appropriate subcategory of
Lemma~\ref{C:relative-Swan-MF-bridge} remains to be proved.
Even a fully faithful embedding of a primary support category
would not by itself prove injectivity on $K_0$ after allowing
all other finite pencil supports.  Neither that injectivity
nor a Swan obstruction is asserted by the proposition.
